\section*{الفصل \ref{ch:g7:circuits-series-parallel} --- الدارات: الحلقات والتسلسل والتفرع}

\begin{solution}{exo:g7:circuits-series-parallel:1}
جريان جسيمات مشحونة تسير في الدارة تحت دفع \omterm{def:g3:first-electric-circuit:battery}{البطارية}. ولا يوجد
إلا ما دامت حلقة مغلقة واحدة على الأقلّ تصل القطبين؛ فافتح
الحلقة يتوقّف المسير في كل مكان دفعة واحدة.
\end{solution}

\begin{solution}{exo:g7:circuits-series-parallel:2}
خارج \omterm{def:g3:first-electric-circuit:battery}{البطارية}، من $+$ حول الدارة إلى $-$. وعلى حلقة \omterm{def:g3:first-electric-circuit:bulb}{المصباح}
الواحد تغادر الأسهم $+$، وتمرّ \omterm{def:g3:first-electric-circuit:bulb}{بالمصباح}، وتعود إلى $-$.
\end{solution}

\begin{solution}{exo:g7:circuits-series-parallel:3}
الخطأ القائل إن العناصر تستهلك التيار، فينبغي أن تتوهّج
المصابيح ``التالية'' أضعف. وتساوي السطوع --- وهو ينجو من كل
تبديل --- يبيّن أن التيار الكامل نفسه يمرّ بكليهما: فلا شيء
يُؤكَل في الطريق.
\end{solution}

\begin{solution}{exo:g7:circuits-series-parallel:4}
هما تيار واحد بعينه: فالحلقة الواحدة تحمل مسيرًا واحدًا عبر كل
عنصر عليها.
\end{solution}

\begin{solution}{exo:g7:circuits-series-parallel:5}
ينقسم --- فيسلك جزء من المسير كل \omterm{def:g5:series-parallel-circuits:parallel}{فرع} --- وعند \omterm{def:g5:series-parallel-circuits:parallel}{العقدة} الثانية
تلتقي الأجزاء كاملةً. وينعم كل \omterm{def:g5:series-parallel-circuits:parallel}{فرع} بدفع \omterm{def:g3:first-electric-circuit:battery}{البطارية} كاملًا.
\end{solution}

\begin{solution}{exo:g7:circuits-series-parallel:6}
وحيدًا $=$ مع جار تفرع واحد (سطوع كامل)، ثم مع رفيق تسلسل واحد
(خافت)، ثم مع رفيقَي تسلسل (أخفت بعد).
\end{solution}

\begin{solution}{exo:g7:circuits-series-parallel:7}
التفرع: ففرعان يسحب كل واحد مسيرًا كاملًا يضاعفان تسليم
\omterm{def:g3:first-electric-circuit:battery}{البطارية} الكلّي، بينما يرسل زوج التسلسل مسيرًا واحدًا مخنوقًا
عبر الاثنين. والسطوع الكامل يُدفَع ثمنه عند \omterm{def:g3:first-electric-circuit:battery}{البطارية}.
\end{solution}

\begin{solution}{exo:g7:circuits-series-parallel:8}
الكراسي هي السائرون المشحونون، يدورون بلا انقطاع ولا يُستهلكون؛
والمتزلّجون هم \omterm{def:g5:everyday-energy:energy}{الطاقة}، تُحمَّل عند \omterm{def:g3:first-electric-circuit:battery}{البطارية} وتُنزَّل عند
المصابيح ضوءًا ودفئًا. والذي ينفد هو مخزون المنتجع من
المتزلّجين --- أي \omterm{def:g5:everyday-energy:energy}{طاقة} \omterm{def:g3:first-electric-circuit:battery}{البطارية} المخزونة.
\end{solution}

\begin{solution}{exo:g7:circuits-series-parallel:9}
كل حلقة تمرّ \omterm{def:g3:first-electric-circuit:bulb}{بالمصباح} A تمرّ \omterm{ex:g3:first-electric-circuit:switch}{بالقاطعة} ثم \omterm{def:g3:first-electric-circuit:bulb}{بالمصباح} B أو
\omterm{def:g3:first-electric-circuit:bulb}{بالمصباح} C --- فرفاق A \omterm{def:g4:series-circuits:series}{على
التسلسل} هما \omterm{ex:g3:first-electric-circuit:switch}{القاطعة} و\emph{زوج} الفرعين ككلّ، لا B وحده ولا C
وحده. وقطع \omterm{def:g5:series-parallel-circuits:parallel}{فرع} C يرسل المسير كله عبر B: فيظلّ A متوهّجًا (على
نحو مختلف قليلًا --- طريق واحد بدل طريقين --- كما ستقيس أجهزة
السنة القادمة).
\end{solution}

\begin{solution}{exo:g7:circuits-series-parallel:10}
قاطعة رئيسية واحدة على الطريق المشترك بعد المغذّي مباشرة، ثم
عشرة فروع متفرّعة، في كل \omterm{def:g5:series-parallel-circuits:parallel}{فرع} \omterm{def:g3:first-electric-circuit:bulb}{مصباح} واحد. فقانون التفرع يمنح كل
\omterm{def:g3:first-electric-circuit:bulb}{مصباح} سطوعًا كاملًا مستقلًّا؛ ووضع \omterm{ex:g3:first-electric-circuit:switch}{القاطعة} الرئيسية \omterm{def:g4:series-circuits:series}{على التسلسل}
--- على الطريق الذي تتقاسمه كل حلقة --- يتيح لقاطعة واحدة أن
تأمر العشرة.
\end{solution}

\begin{solution}{exo:g7:circuits-series-parallel:11}
قسمة \omterm{def:g5:series-parallel-circuits:parallel}{العقدة} ليست قسمة متساوية: فطريق \omterm{def:g5:series-parallel-circuits:parallel}{الفرع} نفسه قد يقاوم
المسير، والسلك الطويل الرفيع يخنق حصة فرعه --- سائرون أقلّ في
الثانية، وتوهّج أضعف. فالطرق نفسها تقاوم التيار: وهي الفكرة
وراء مقاومة السنة القادمة.
\end{solution}

\begin{solution}{exo:g7:circuits-series-parallel:12}
X، وحده على فرعه، ينال الدفع كاملًا: فهو الأسطع. وY وZ يتقاسمان
دفع فرعهما \omterm{def:g4:series-circuits:series}{على التسلسل}: فهما خافتان بالتساوي، وكلاهما دون X.
فإن احترق Z انقطعت حلقته: فيعتم Y معه --- ويواصل X سطوعه غير
مبالٍ. وإن احترق X مات \omterm{def:g5:series-parallel-circuits:parallel}{فرع} X وحده: فيواصل Y وZ، خافتين كما
كانا.
\end{solution}

\begin{solution}{pb:g7:circuits-series-parallel:1}
\textbf{1.} أربعة فروع متفرّعة، واجب في كل واحد: \omterm{def:g5:series-parallel-circuits:parallel}{فعلى التفرع}
يعمل كل \omterm{def:g5:series-parallel-circuits:parallel}{فرع} بكامل قوّته مهما فعل غيره. وأيّ اقتران \omterm{def:g4:series-circuits:series}{على التسلسل}
كان سيتيح لقاطعة \omterm{def:g3:first-electric-circuit:bulb}{مصباح} قراءة أو \omterm{def:g3:first-electric-circuit:bulb}{لمصباح} درج محترق أن يُخفت
\omterm{def:g3:first-electric-circuit:bulb}{المصباح} العظيم أو يُعتمه --- وذلك لا يُغتفر.
\textbf{2.} \omterm{ex:g3:first-electric-circuit:switch}{القاطعة} الرئيسية على الطريق المشترك بعد المجموعة
مباشرة؛ وكل واحدة من قواطع الفروع الأربع \omterm{def:g4:series-circuits:series}{على التسلسل} مع واجبها،
داخل فرعها.
\textbf{3.} نعم --- بقانون التفرع: فالفراغ في \omterm{def:g5:series-parallel-circuits:parallel}{فرع} لا يوقف إلا
حصة ذلك \omterm{def:g5:series-parallel-circuits:parallel}{الفرع}؛ وسائر الفروع تواصل مسيرها.
\textbf{4.} المجموعة، \omterm{ex:g3:first-electric-circuit:switch}{فالقاطعة} الرئيسية، ثم \omterm{def:g5:series-parallel-circuits:parallel}{عقدة} تتشعّب إلى
أربعة فروع --- (\omterm{def:g3:first-electric-circuit:bulb}{المصباح} العظيم $+$ قاطعته)، و(بوق الضباب $+$
قاطعته)، و(\omterm{def:g3:first-electric-circuit:bulb}{مصباح} القراءة $+$ قاطعته)، و(\omterm{def:g3:first-electric-circuit:bulb}{مصباح} الدرج $+$
قاطعته) --- ثم تلتقي عائدة إلى المجموعة.
\textbf{5.} نعم: فكل \omterm{def:g5:series-parallel-circuits:parallel}{فرع} يمتدّ بين جانبَي المجموعة، فينال
الدفع كاملًا مهما كثر الجيران الشاربون إلى جانبه.
\textbf{6.} لا: فالتيار المارّ \omterm{def:g3:first-electric-circuit:bulb}{بالمصباح} يغادره كاملًا ---
فالكراسي تدور غير مأكولة. وإنما الذي يأخذه \omterm{def:g3:first-electric-circuit:bulb}{المصباح} \omterm{def:g5:everyday-energy:energy}{طاقة} من
مخزن المجموعة، يسلّمها المسير العابر.
\textbf{7.} أربعة مسايير كاملة دفعة واحدة: \omterm{def:g5:everyday-energy:energy}{فطاقة} المجموعة
تُستنزف بأربعة واجبات سريعًا. والفاتورة تُدفَع في غرفة
البطاريات --- من \omterm{def:g5:everyday-energy:energy}{الطاقة} المخزونة، لا من تيار ضائع أبدًا.
\textbf{8.} من القطب $+$ في المجموعة، عبر \omterm{ex:g3:first-electric-circuit:switch}{القاطعة} الرئيسية،
إلى \omterm{def:g5:series-parallel-circuits:parallel}{العقدة}؛ ثم إلى \omterm{def:g5:series-parallel-circuits:parallel}{فرع} \omterm{def:g3:first-electric-circuit:bulb}{مصباح} القراءة، عبر قاطعته المغلقة
\omterm{def:g3:first-electric-circuit:bulb}{والمصباح}؛ ثم عبر \omterm{def:g5:series-parallel-circuits:parallel}{العقدة} حيث تلتقي الفروع الثلاثة الأخرى
(المعتمة)؛ ثم عودة إلى القطب $-$. وهو يتجنّب فروع \omterm{def:g3:first-electric-circuit:bulb}{المصباح}
العظيم وبوق الضباب والدرج كلها.
\textbf{9.} لا شيء غيره --- فلا يموت إلا \omterm{def:g5:series-parallel-circuits:parallel}{فرع} الدرج مع مصباحه:
وهو استقلال قانون التفرع.
\textbf{10.} \omterm{def:g6:circuits-and-safety:short}{دارة قصيرة} قبل كل \omterm{def:g5:series-parallel-circuits:parallel}{فرع}. فيهجر المسير كل طريق عامل
إلى الاختصار العاري بلا جهد: فتجوع المصابيح والبوق بينما يحمل
سلك الاختصار اندفاعًا جامحًا ساخنًا --- ويفشل كل واجب دفعة
واحدة حتى يُنزع السلك.
\textbf{11.} امنح \omterm{def:g3:first-electric-circuit:bulb}{المصباح} العظيم حلقته المستقلّة على المجموعة
الثانية: مجموعة ثانية، وقاطعتها الخاصة، \omterm{def:g3:first-electric-circuit:bulb}{والمصباح} العظيم ---
لا يتقاسم سلكًا مع اللوحة العامة. فيُسكت اختصار عبر المجموعة
الأولى كل شيء \emph{إلا} \omterm{def:g1:light-and-shadows:source}{ضوء} السفن.
\textbf{12.} مثلًا: ``حفظ قانون التفرع سيادة كل واجب الليلة.
ودحضت الليلة، مرة أخرى، الاعتقاد بأن المصابيح تلتهم التيار ---
فقد عاد المسير كاملًا وإن استُنزفت المجموعات. وسلك عارٍ واحد،
اختصر المجموعة، كلّفنا أسوأ ساعة: كل طريق خالٍ والاختصار
يحترق.''
\end{solution}
